% `tex.web` §863 and §876: which comes first in the log, the reasoning or the
% complaint.
%
% `line_break` writes §846's narration as each pass runs, and only then does
% §876's `post_line_break` hand each line to `hpack` — which is where an
% underfull or overfull line is reported (§660/§663). So the trace of how a
% paragraph was broken always precedes the warning about a line it chose.
%
% Ours drained the narration after the warnings, so a paragraph that went loose
% announced the loose line above the reasoning that produced it. In trip.log
% that reads as a missing trace and a misplaced warning at once, which is two
% wrong diagnoses of one fault.
%
% The lines here are deliberately loose rather than overfull: an overfull line
% means no feasible break, which sends TeX into a second pass whose narration
% needs §851's artificial demerits to match, and that is its own piece of work.
\catcode`\{=1 \catcode`\}=2
\tracingonline=1
\tracingparagraphs=1
\hbadness=0
\hfuzz=0pt
% High enough that the first pass accepts the loose line: what the *second*
% pass narrates needs §851's artificial demerits to match, and that is its own
% piece of work.
\pretolerance=10000
\showboxbreadth=2 \showboxdepth=1
\hsize=100pt \parindent=0pt
\font\rip=cmr10 \rip

% Short enough that the interword glue must stretch a long way to justify it,
% which is what makes the line loose rather than the break wrong.
AAA AAA\par
AAA AAA\par
\end
